\documentclass[../../../include/open-logic-section]{subfiles}

\begin{document}

\olfileid{sth}{choice}{justifications}
\olsection{চয়নের অন্তর্নিহিত ন্যায্যতা}

বৃহত্তর প্রশ্ন হলো: সুক্রমবিন্যাস, চয়ন, কিংবা
আকারের বিচারে সব সেটের তুলনীয়তা—যেটিই ধরি—
তার ন্যায্যতা কি দেখানো যায়?

একটি \emph{অন্তর্নিহিত} ন্যায্যতার চেষ্টা দেখি।
\olref[z][story]{sec}-এ স্তরক্রম সম্বন্ধে
কয়েকটি নীতি এনেছিলাম। তার একটি আবার বলা দরকার:
\begin{enumerate}
	\item[] \stagesacc। যেকোনো পর্যায় $S$-তে,
	$S$-এর \emph{আগে} গঠিত যেকোনো সেটগুলি
	নিলে, $S$-তে ঠিক সেগুলিকেই সদস্য করে একটি
	সেট গঠিত হয়। ওই পর্যায়ে অন্য কিছু গঠিত হয় না।
\end{enumerate}
বস্তুত বহু লেখক বলেছেন, এই বা এর কাছাকাছি
এক নীতি দিয়ে চয়নের স্বতঃসিদ্ধের ন্যায্যতা
দেখানো যায়। পদ্ধতিটির একটি সংক্ষিপ্ত ব্যাখ্যা দেব।

প্রথমে একটি ছোট ফল দেখি। এটি চয়নের \emph{আরও
একটি} সমতুল্য রূপ দেয়:

\begin{thm}[$\ZF$-এ]\ollabel{choiceset}
চয়ন নিচের নীতিটির সমতুল্য। $A$-র !!{element}s
পরস্পর বিচ্ছিন্ন এবং অশূন্য হলে এমন $C$ আছে
যে প্রত্যেক $x\in A$-র জন্য $C\cap x$
একপদী সেট। (এমন $C$-কে $A$-র
{চয়ন-সেট} বলি।)
\end{thm}

ফলটির প্রমাণ সরল; পাঠকের অনুশীলন হিসেবে রইল।

\begin{prob}
\olref[sth][choice][justifications]{choiceset}
প্রমাণ করো। অসুবিধা হলে
\cite[pp.~242--3]{Potter2004}-এ প্রমাণ পাবে।
\end{prob}

মূল কথা, $A$-র চয়ন-সেটকে সংযোগের মধ্যে সীমিত করলে
তা $A$-র কোনো চয়ন-অপেক্ষকের বিস্তার। তাই চয়নের ন্যায্যতা
দেখাতে চাইলে সমতুল্য চয়ন-সেটের অস্তিত্বের
ন্যায্যতা দেখার চেষ্টা করতে পারি। এবার সেটিই করব।
% BN-SRC-828 TARGET-CORRECTION-BEGIN
\par\noindent\textbf{পাঠ-সংশোধন।} উপপাদ্যের সংজ্ঞায় $C$-র মধ্যে $\bigcup A$-র বাইরের
অপ্রাসঙ্গিক সদস্যও থাকতে পারে। তাই আগে
$C'=C\cap\bigcup A$ নিতে হয়; তখন প্রত্যেক $x\in A$-র
সঙ্গে তার একমাত্র ছেদের সদস্যকে $f(x)$ করলে
$\ran{f}=C'$। উল্টো দিকে চয়ন-অপেক্ষকের বিস্তার প্রত্যেক
বিচ্ছিন্ন সদস্য-সেটকে ঠিক একবার ছেদ করে। arbitrary অশূন্য
সদস্য-সেটের পরিবারে সমতুল্যতার প্রমাণের জন্য আগে
$x$-কে $\{x\}\times x$ দিয়ে চিহ্নিত করি। এই পরিবার
বিচ্ছিন্ন; তার চয়ন-সেটের অনন্য সদস্যের দ্বিতীয় স্থানাঙ্ক
মূল $x$ থেকে চয়ন দেয়। শূন্য সদস্য-সেট বাদ যায়; শূন্য
পরিবারের সাক্ষী শূন্য অপেক্ষক।
% BN-SRC-828 TARGET-CORRECTION-END

$A$-র !!{element}s পরস্পর বিচ্ছিন্ন ও অশূন্য
ধরি। \stageshier{} অনুযায়ী (দেখো
\olref[z][story]{sec}) $A$ কোনো পর্যায়~$S$-তে
গঠিত। লক্ষ করো, $\bigcup A$-র সব !!{element}s
পর্যায়~$S$-এর আগেই পাওয়া যায়। এবার
\stagesacc{} বলে, $S$-এর আগে গঠিত
\emph{যেকোনো} সেটগুলিকে ঠিক সদস্য করে
একটি সেট $S$-তে গঠিত হয়। অন্যভাবে বললে,
আগেই পাওয়া সেটগুলির প্রত্যেক \emph{সম্ভাব্য}
সংকলন $S$-তে আছে। কিন্তু $A$-র চয়ন-সেট
গড়ার মতো বস্তু বেছে নেওয়া নিশ্চয়ই
\emph{সম্ভব}; সেটি $\bigcup A$-র একটি
নির্দিষ্ট উপসেট। অতএব প্রয়োজনমতো এমন
একটি চয়ন-সেট আছে।
% BN-SRC-829 TARGET-CORRECTION-BEGIN
\par\noindent\textbf{পাঠ-সংশোধন।} এটি অন্তর্নিহিত ন্যায্যতার একটি দার্শনিক প্রস্তাব, চয়নকে
$\ZF$ থেকে প্রমাণ করার formal ধাপ নয়। আগে থাকা সেটগুলির
সব উপসেট গড়তে পারা এবং প্রত্যেক পরিবারের জন্য এক সদস্যের
নির্বাচনসমষ্টির অস্তিত্ব এক কথা নয়। ``এমনভাবে বেছে নেওয়া
সম্ভব''-এর সেটতাত্ত্বিক অর্থেই একটি চয়ন-সেটের অস্তিত্ব
প্রবেশ করে, আর উপপাদ্য অনুসারে সেই অস্তিত্ব চয়নের সমতুল্য।
এই নতুন অনুমানকে কেবল ঘাতসেট বা পৃথকীকরণ বলে লুকানো হয়নি।
% BN-SRC-829 TARGET-CORRECTION-END

অন্তর্নিহিত ভিত্তিতে চয়নের ন্যায্যতা দেখানোর
এটি খুব সংক্ষিপ্ত চেষ্টা। আরও এগোতে চাইলে
পটারের সুন্দর বিস্তার
(\citeyear[\S14.8]{Potter2004}) পড়তে পারো।

\end{document}
